problem:  
Let \( (X_t, \omega_t, \Omega_t) \) be a smooth 1‑parameter family of compact Calabi–Yau \(n\)-manifolds equipped with their Ricci‑flat Kähler metrics \( g_t \) lying in fixed cohomology classes.  
Assume that as \(t \to 0\), the metrics \(g_t\) *collapse with bounded curvature away from a singular set*, in the sense that  
\[
(X_t, g_t) \longrightarrow (B, g_B)
\]
in the Gromov–Hausdorff topology, where \(B\) is a real \(n\)-dimensional manifold endowed with a Riemannian metric \(g_B\) of Monge–Ampère type.  
Denote by \(C_{p,t}\) the tangent cone at a point \(p \in X_t\) with respect to the rescaled metric \(r^{-2}g_t\) as \(r \to 0\).

**Question:**  
Can one determine, from the asymptotic geometry of the family of tangent cones \(\{ C_{p,t} \}_{p \in X_t,\ t \to 0}\), whether the Gromov–Hausdorff limit \(B\) admits a *smooth integral affine structure* compatible with a (possibly singular) special Lagrangian torus fibration predicted by the Strominger–Yau–Zaslow (SYZ) conjecture?  
Equivalently: does the local metric structure encoded in the tangent cones suffice to detect the emergence of affine coordinates on \(B\) associated with families of special Lagrangian tori in \(X_t\)?

Why is it a "good" problem:  
This question lies precisely at the uncharted intersection of **metric limit theory** (Cheeger–Colding–Tian analysis of tangent cones in collapsing Ricci‑flat spaces) and **mirror‑symmetric geometry** (the SYZ mechanism producing real affine bases from Calabi–Yau collapse). Current theory describes how such collapses lead to Monge–Ampère metrics on \(B\), but no known framework links the **microscopic tangent cone geometry** to the **emergence of global integral affine structures** governing mirror duality.  

It is mathematically sound, conceptually simple to state yet profoundly deep—asking whether infinitesimal metric data alone encode the torus‑symmetry and local potential structure of SYZ bases. A resolution would unify analytic notions of Ricci‑flat collapse with the symplectic/topological data of mirror symmetry, creating a bridge between differential geometry, geometric analysis, and complex/symplectic topology.  
Because it seeks local‑to‑global reconstruction of one of the central conjectural structures in Calabi–Yau geometry via tools not yet known to interact, it is an exemplary “good” problem: precise, minimal, and fundamentally transformative in scope.